% ============================================================================
%  front/intro.tex — مقدمه، راهنمای استفاده، نمادها، توصیهٔ مطالعه
% ============================================================================

% ------------------------------------------------------------- مقدمه --------
\JFrontSec{مقدمه}
\begin{JNass}[اَلسَّلامُ عَلَيْكُم \,\textbar\, با سلام]
عزیزِ کوشا، به جزوهٔ «عربی، زبانِ قرآن» پایهٔ دوازدهم خوش آمدی. آنچه در دست
داری، حاصلِ گردآوری و بازآراییِ متنِ کتابِ درسی، ترجمه‌ها، واژه‌نامه‌ها و
قواعدِ هر پنج درس است؛ همراه با خودآزمایی‌ها و تمرین‌های کتاب (بدونِ پاسخ،
برای تمرینِ ذهنی)، پاسخ‌نامهٔ تشریحیِ آن‌ها در پایانِ هر درس، بانکِ تست‌های
اضافی و آزمون‌های جامعِ مشابهِ کنکور و در پایانِ کار، واژه‌نامهٔ جامعِ
الفبایی و درسنامهٔ کاملِ قواعد برای مرورِ شبِ امتحان.
\end{JNass}

چرا عربی مهم است؟ زیرا این درس ترکیبی از \textbf{واژگان}، \textbf{ترجمه} و
\textbf{قواعد} است و در کنکورِ علومِ انسانی سهمِ پررنگی دارد. خبرِ خوب این‌که
عربیِ کنکور «قاعده‌محور» است: بخشِ بزرگی از تست‌ها با تسلط بر همین قواعدی که
در این جزوه درس‌به‌درس آموخته‌ای حل می‌شوند. پس با برنامه پیش برو، هر درس را
در سه مرحله بخوان و از پاسخ‌نامه‌ها فقط برای «آزمودنِ» خودت استفاده کن.

% --------------------------------------------------------- راهنمای استفاده --
\JFrontSec{راهنمایِ استفاده از جزوه}
هر درسِ این جزوه یک ساختارِ ثابت دارد:

\begin{itemize}
  \item \textbf{آغازینِ درس:} شمارهٔ درس در مدالیونِ مرکزی، عنوانِ عربی و فارسی،
        و حکمتِ آغازین با ترجمه.
  \item \textbf{متنِ درس (اَلنَّصُّ):} عینِ متنِ کتاب؛ شعرها به‌صورتِ مصراع‌بندی‌شده
        و حاشیه‌نویسی‌های کتاب، درونِ متن آمده است.
  \item \textbf{ترجمهٔ روان (اَلتَّرْجَمَةُ):} برگردانِ فارسیِ کاملِ متن، بیت‌به‌بیت
        و بندبه‌بند.
  \item \textbf{واژه‌نامه (اَلْمُعْجَمُ):} همهٔ واژه‌های مُعجَمِ کتاب، به‌همراهِ
        روابطِ معنایی (مترادف، متضاد، جمع و مفرد) و یادداشت‌ها.
  \item \textbf{قواعد (إِعْلَمُوا):} درسنامهٔ کاملِ قواعدِ همان درس، با جدول‌ها،
        مثال‌های متن و نکته‌های کنکوری.
  \item \textbf{خودآزمایی و تمرین‌ها:} عینِ پرسش‌های کتاب، \emph{بدونِ پاسخ}؛
        این بخش «میدانِ تمرینِ» توست.
  \item \textbf{پاسخ‌نامه:} در پایانِ هر درس، پاسخِ تشریحیِ خودآزمایی‌ها و
        تمرین‌های کتاب — پیش از دیدنِ آن‌ها، خودت را بیازما!
\end{itemize}

\begin{JGood}[توصیهٔ مطالعه]
\textbf{نوبتِ نخست:} متن و ترجمه را همزمان بخوان و واژه‌های ناآشنا را نشانه بگیر.
\textbf{نوبتِ دوم:} قواعدِ درس را با مثال‌های کتاب تمرین کن. \textbf{نوبتِ سوم:}
خودآزمایی و تمرین‌ها را پاسخ بده و با پاسخ‌نامه بسنج. در هفتهٔ پیش از آزمون نیز
«مرورِ سریع و درسنامهٔ جامعِ قواعد» در انتهای جزوه برای بازبینیِ کاملِ دروس به
کار تو می‌آید.
\end{JGood}

% --------------------------------------------------------------- نمادها -----
\JFrontSec{راهنمایِ نمادها و جعبه‌ها}
\begin{JLegend}
\begin{itemize}
  \item \textbf{متنِ درس (JNass):} عینِ متنِ درس، روی زمینهٔ کاغذیِ گرم.
  \item \textbf{ترجمهٔ روان (JTarjome):} برگردانِ فارسیِ متن، بیت‌به‌بیت.
  \item \textbf{واژه‌نامه (JMoajam):} جدولِ واژه‌های مُعجَمِ درس.
  \item \textbf{قاعده (JQaede):} آموزشِ کاملِ قواعد — نیلی.
  \item \textbf{نکته (JNokte):} نکته‌های مهم و کنکوری.
  \item \textbf{دامِ آزمون (JDam):} خطاهای رایج — قرمز.
  \item \textbf{خوب است بدانیم (JGood):} فراتر از کتاب.
  \item \textbf{خودآزمایی (JTest):} پرسش‌های خودآزماییِ کتاب، بدونِ پاسخ — ارغوانی.
  \item \textbf{تمرین (JTamrin):} تمرین‌های کتاب، بدونِ پاسخ — نیلیِ روشن.
  \item \textbf{پاسخ‌نامه (JKey):} پاسخِ تشریحی — سبز.
\end{itemize}
\end{JLegend}
